\section{Four rows}\label{sec:four}

\begin{theorem}[four rows]\label{thm:fourvalue}
For every $m\ge6$ we have $f(4,m)=m+4$.
\end{theorem}

To attain the upper bound of Theorem~\ref{thm:rowbound}, use
\begin{equation}\label{eq:fourS}
S=\{(i,1):2\le i\le m-1\}\cup\{(m-1,2),(m,2),(1,3),(2,3),(1,4),(m,4)\}.
\end{equation}
Here $|S|=m+4$ and $G_S$ is the eight-cycle $r_1\,c_2\,r_3\,c_1\,r_4\,c_m\,r_2\,c_{m-1}$ with leaves $c_3,\ldots,c_{m-2}$ on $r_1$.

\begin{proof}[Proof of Theorem~\ref{thm:fourvalue}]
The graph $G_S$ has no four- or six-cycle. Thus Lemma~\ref{lem:girth}
distinguishes the axis-aligned rectangles, and none is fully hidden.

For tilted rectangles, use Lemma~\ref{lem:tilted-narrow}.
Their corners in row $1$ are black and those in row $4$ are white.
For $D(i,1)$ the additional black corners are $(i+2,2)$ when $i\ge m-3$
and $(i+1,3)$ when $i=1$. For $D(i,2)$ they are $(i,3)$ when $i\le2$
and $(i+1,2)$ when $i=m-2\ge4$. For $\kA_i$ and $\kB_i$ they are
$(i,3)$ when $i\le2$ and $(i+3,2)$ when $i\ge m-4$.
Three shown corners suffice, leaving the weak rectangles. Below, $s=1$ for $\kA_i$
and $s=2$ for $\kB_i$.
\begin{center}\small
\begin{tabular}{@{}lll@{}}
\toprule
rectangle & which $i$ & corners shown\\
\midrule
$D(1,1)$ & & $(1,2),(3,2)$\\
$D(i,1)$ & $i\in\{m-3,m-2\}$ & $(i,2),(i+1,3)$\\
$\kA_i$ or $\kB_i$ & $i\in\{1,2\}$, $i\notin\{m-4,m-3\}$ & $(i+s,4),(i+3,2)$\\
$\kA_i$ or $\kB_i$ & $i\in\{m-4,m-3\}$, $i\notin\{1,2\}$ & $(i,3),(i+s,4)$\\
$\kA_i$ or $\kB_i$ & $i=2$, $m=6$ & $(i+s,4)$\\
\bottomrule
\end{tabular}
\end{center}
The first four cases have different shown row sets. Within any one,
the point in row $2$ or $3$ determines $i$, and the row-$4$ point
determines $s$.
For rivals, no row $r$ has
both $(1,r)$ and $(3,r)$ black, ruling out a rival of
$\{(1,2),(3,2)\}$. In the next three cases, the axis-aligned rectangle
through the pair has a white corner, respectively $(i,3)$, $(i+s,2)$,
or $(i,4)$. Here we use $i\ge3$, $i+s\le4<m-1$, and
$3\le i\le m-3$, respectively. The final case has $m=6$ and one
shown point $(z,4)$ with $z\in\{3,4\}$. Column $z$ is black only in
row $1$, and no column is black in rows $1$ and $4$, so this singleton
also has no rival.
\end{proof}

For $m\ge4$, Theorem~\ref{thm:equality} leaves
$3\,(m)_4\,4^{m-4}$ \emph{candidates} of size $m+4$.
Each has four doubles forming a four-cycle $F$ on the rows,
with a singleton in every remaining column.

\begin{corollary}[widths four and five, computer-assisted]\label{cor:fourboundary}
We have $f(4,4)=7$ and $f(4,5)=8$.
\end{corollary}

\begin{proof}
The boundary cross gives the lower bounds. For the upper bounds,
the candidate search finds an ambiguous pair in every one of the $72$
and $1{,}440$ cases, respectively.
\end{proof}

\begin{theorem}[no hidden rectangle]\label{thm:hidden}
For $m\ge4$, no valid blackout of the $4\times m$ grid hides a rectangle.
\end{theorem}

\begin{proof}
By Lemma~\ref{lem:diff}, it is enough to express each rectangle's
corner set as $C(R_1)\triangle C(R_2)$ for $R_1\ne R_2$.
For an axis-aligned rectangle, use its two rows and pair each column
$x,y$ with a third column $w$.
For $1\le i\le m-3$, put
$\mathsf Q_i=\{(i,2),(i,3),(i+3,2),(i+3,3)\}$. The tilted cases follow from
\begin{equation}\label{eq:fourxor}
\begin{gathered}
\kC_i=\kA_i\,\triangle\,\mathsf Q_i,\qquad \kE_i=\kB_i\,\triangle\,\mathsf Q_i,\\
D(i+1,1)=D(i,2)\,\triangle\,\kA_i,\qquad D(i+1,2)=D(i,1)\,\triangle\,\kE_i,
\end{gathered}
\end{equation}
since any of the three sets in an identity is the difference of the
other two. The first line handles $\kA_i$, $\kB_i$, $\kC_i$, and $\kE_i$.
For a diamond $D(j,1)$ or $D(j,2)$, use the second line with $i=j-1$
or $i=j$. One choice is always available, since $j\ge2$ or
$j\le m-3$ when $m\ge4$.
\end{proof}

\begin{theorem}[forbidden-block bound]\label{thm:foursurvival}
For every $m\ge6$ we have $a_4(m)\le3\,(m)_4\,4^{m-4}(53/64)^k$, where
$k=\max(0,\lfloor m/4\rfloor-4)$. So $a_4(m)=O\bigl(m^4\,212^{m/4}\bigr)$.
\end{theorem}

\begin{proof}
First choose $F$, the four double columns, and their pairs, in
$3\,(m)_4$ ways. Call four consecutive singleton columns a
\emph{block}, recorded by their black rows. Starting at column $x$,
the four \emph{template words} $2413$, $3412$, $2143$, and $3142$
are the corner rows of $\kC_x$, $\kA_x$, $\kE_x$, and $\kB_x$.
Fix one template rectangle $T$ with word $t_1t_2t_3t_4$.

Swapping $t_p$ and $t_q$ gives a \emph{swap word}. It cannot occur,
since $T$ and the axis-aligned rectangle through $(x+p-1,t_p)$ and
$(x+q-1,t_q)$ would show the same pair. The templates are exactly the
permutations whose end letters both lie in $\{2,3\}$. Swaps within
the end or middle positions give templates again. Swaps between these
positions give the sixteen permutations with one end letter in
$\{2,3\}$. Together these are $20$ forbidden swap words.

An \emph{edit word} changes one letter $t_p=r$ to $s$ with
$\{r,s\}\in F$. This is also impossible. The rectangle $T$ and the
axis-aligned rectangle on rows $r,s$ and columns $x+p-1$ and $d$,
where $d$ is the double column for $\{r,s\}$, would show only
$(x+p-1,r)$. There are $8$ edits per template. Two templates can
share an edit only when two positions contain $r,s$ in opposite orders.
They then share two edits if $\{r,s\}\in F$. This happens for two
template pairs with $\{r,s\}=\{2,3\}$ and two with $\{1,4\}$.
Hence at least $32-8=24$ edit words are forbidden. They repeat a
letter and so are separate from the swap words. A block therefore
allows at most $256-20-24=212=4^4\cdot53/64$ words.

Partition the columns into the $\lfloor m/4\rfloor$ groups
$4j-3,\ldots,4j$. At most four contain a double, leaving at least
$k$ singleton blocks. All other singleton columns have at most $4$
choices, which gives the bound.
\end{proof}

\begin{lemma}[six columns of memory]\label{lem:fourmemory}
Let $m\ge6$, and let $F$ be the edge set of a four-cycle on the rows. A word over the
four singleton letters and the pairs of $F$, each pair used once, represents an optimum
if and only if (i) no tilted rectangle is hidden, (ii) no tilted rectangle shares its
presentation with an axis-aligned rectangle whose columns are among its own and the
double columns, and (iii) no two distinct tilted rectangles in a block of at most seven
consecutive columns share a presentation.
\end{lemma}

\begin{proof}
Necessity follows from validity and Theorem~\ref{thm:hidden}.
Conversely, the prescribed doubles and singletons give an axis-valid
word by Theorem~\ref{thm:equality}. No axis-aligned rectangle is hidden,
since that would give a four-cycle in $G_S$, and (i) excludes hidden
tilted rectangles.

If two rectangles nevertheless show the same $P$, then
$P\ne\emptyset$ and at least one is tilted. If the other is
axis-aligned, a column of it either meets $P$ or has two black corners.
In the first case it is a column of the tilted rectangle, and in the
second it is a double column. Thus (ii) excludes the pair.
If both are tilted, choose a common column $c$ meeting $P$.
Each rectangle spans at most four columns, so their union fits in
$c-3$ to $c+3$, contradicting (iii). The word is valid and has
$m+4$ black points, hence is an optimum.
\end{proof}

For a word built from singleton letters and pairs of $F$, a
tilted rectangle \emph{fails locally} in any of three cases. It is
hidden, its presentation has a rival under Lemma~\ref{lem:rival}
with pair set $F$, or another tilted rectangle shows the same set.

\begin{proposition}[the automaton]\label{prop:fourautomaton}
Let $F$ be the edge set of a four-cycle on the rows. The automaton $\mathcal A_F$ reads
the four singleton letters and the pairs in $F$. It starts at $(\epsilon,\emptyset)$,
where $\epsilon$ is the empty word. A state $(u,U)$ records the suffix of at most six
letters read so far and the set $U\subseteq F$ of pairs already used.

On reading $y$, it rejects if $y\in U$ or if a tilted rectangle in $uy$ ending in the
new column fails locally. Otherwise it replaces $u$ by the suffix of $uy$ of length
$\min(6,|uy|)$ and $U$ by $U\cup(\{y\}\cap F)$. A word is accepted if it is read to the end
without rejection and the final register is $U=F$. For each $m\ge6$, the accepted words
of length $m$ are exactly the optima of the $4\times m$ grid with pair set $F$.
\end{proposition}

\begin{proof}
Take $m\ge6$. At its last column, each tilted rectangle lies in $uy$, since its span is at most three.
An accepted word uses every pair of $F$ once, so it is a candidate.
Presentations of at least three points are already unique
(Lemma~\ref{lem:threecorners}). The local tests rule out hidden
tilted rectangles and axis-aligned rivals, giving (i) and (ii) of
Lemma~\ref{lem:fourmemory}. For (iii), a pair of tilted rectangles
within seven columns is present when the later last column is read.
That step compares their presentations. The memory lemma therefore
makes an accepted word an optimum.

In the other direction, an optimum with pair set $F$ uses every pair
once. Theorem~\ref{thm:hidden} excludes hidden rectangles, and validity
excludes every repeated presentation. Applied with $F$,
Lemma~\ref{lem:rival} cannot find a rival. The word passes each step
and finishes with register $F$, so it is accepted.
\end{proof}

The script \nolinkurl{four_row_equivalence_check.py} builds $\mathcal A_F$
($60{,}975$, $138{,}095$, and $55{,}575$ states) and checks equivalence with
a smaller machine from an earlier search. Including the pair register,
that machine has $2{,}519$, $3{,}089$, and $2{,}119$ reachable states.
Two transfer counters agree for $m\le1500$, and direct candidate tests agree
for $m\le10$ (Table~\ref{tab:fourvalues}).

\begin{table}[!htb]
\centering\scriptsize\setlength{\tabcolsep}{4pt}
\begin{tabular}{@{}lrrrrrrr@{}}
\toprule
$m$ & $6$ & $7$ & $8$ & $9$ & $10$ & $11$ & $12$\\
\midrule
$\{1,2\}\mid\{3,4\}$ & $0$ & $4$ & $128$ & $920$ & $4{,}852$ & $24{,}112$ & $114{,}424$\\
$\{1,3\}\mid\{2,4\}$ & $20$ & $256$ & $1{,}780$ & $11{,}128$ & $63{,}540$ & $346{,}320$ & $1{,}785{,}768$\\
$\{1,4\}\mid\{2,3\}$ & $6$ & $32$ & $210$ & $1{,}284$ & $7{,}006$ & $32{,}224$ & $137{,}656$\\
$a_4(m)$ & $26$ & $292$ & $2{,}118$ & $13{,}332$ & $75{,}398$ & $402{,}656$ & $2{,}037{,}848$\\
\bottomrule
\end{tabular}
\caption{Four-row optima by pair set (computer-assisted).}
\label{tab:fourvalues}
\end{table}

\begin{theorem}[the four-row generating function, computer-assisted]\label{thm:fourgfexact}
There are coprime integer polynomials $P_4=26z^6-826z^7+\cdots$ of degree $444$ and $Q_4=\sum_{j=0}^{420}q_jz^j$ of degree $420$ with $q_0=1$ and $\sum_{m\ge6}a_4(m)z^m=P_4/Q_4$. The recurrence $\sum_jq_j\,a_4(n-j)=0$ holds for every $n\ge445$ and fails at $n=444$. Over the integers
\[
Q_4=(1-z)^2(1+z)^2\,g_5^3\,\tilde g_5^3\,g_6^3\,g_{11}^5\,g_{17}^3\,g_{19}^3\,g_{20}^5\,g_{21}^5,
\]
where $g_d$ and $\tilde g_d$ are irreducible of degree $d$ with constant term $1$.
\end{theorem}

\begin{proof}
The automaton proposition identifies the counted words with the
optima. Running Berlekamp--Massey modulo eight primes on $a_4(m)$ for
$m\le1500$ produces $P_4$ and $Q_4$.
To check all widths, let $T_F$ be the transfer
matrix on reachable states and write $v(n)=v(0)T_F^n$.
The state vectors satisfy $\sum_jq_jv(n_F-j)=0$ at
$n_F=436$, $435$, and $445$ for the three machines.
Powers of $T_F$ proves the recurrence for
all $n\ge445$. Since $\deg P_4=444$, it fails at $n=444$.
The displayed factorization is obtained in SymPy.
\end{proof}

\begin{corollary}[growth, computer-assisted]\label{cor:fourgrowth}
Let $\rho_4=3.351996777310\ldots$ be the largest real root of $p(x)=x^{20}g_{20}(1/x)=x^{20}-4x^{19}+2x^{18}+\cdots-26x^2-4x+8$. There is a polynomial $B$ of degree $4$ with $a_4(m)=\rho_4^m B(m)+O((1000/341)^m)$, so $a_4(m)\sim c\,m^4\rho_4^m$ with $c\approx8.1611\times10^{-5}$.
\end{corollary}

\begin{proof}
Isolate the roots of every factor of $Q_4$ exactly. The only root
of modulus at most $0.341$ is $r=1/\rho_4=0.29832\ldots$,
simple in $g_{20}$ and therefore of order five in $Q_4$.
Separate its principal part as
$P_4/Q_4=\sum_{j=1}^{5}c_j(1-z/r)^{-j}+E(z)$.
The coefficients of $E$ are $O((1000/341)^m)$.
Coprimality gives $c_5=H(r)\ne0$, where
$H(z)=P_4(z)(1-z/r)^5/Q_4(z)$.
Now $[z^m](1-z/r)^{-j}=\rho_4^m\binom{m+j-1}{j-1}$, so
$B(m)=\sum_jc_j\binom{m+j-1}{j-1}$ and $c=c_5/24$.
Interval arithmetic encloses $c$ between
$8.1610970775\times10^{-5}$ and $8.1610970776\times10^{-5}$.
\end{proof}
